\documentclass[11pt]{article}
\usepackage[utf8]{inputenc}
\usepackage[margin=1in]{geometry}
\usepackage{hyperref}
\usepackage{enumitem}
\usepackage{parskip}

\title{\textbf{Did the AI Really \emph{Need} the Rulebook, or Did It Already Know the Answer?}\\[4pt]
\large A plain-language companion to the paper}
\author{}
\date{}

\begin{document}
\maketitle

\begin{quote}
\small\emph{This is a high-school-level retelling of the technical paper, written so we can
check that the argument actually holds together. It is a companion, not the paper of record;
every claim here is stated more carefully (and hedged) in \texttt{main.tex}. It is a
\textbf{test bench}---a way to check the machinery in simulation---not a proof that real
students learn more.}
\end{quote}

\section*{The problem in one paragraph}

Picture an AI that is supposed to answer using an approved rulebook or reference corpus---say,
the rules that stop ships from colliding, or a set of facts about your company's products. It
gives a right-looking answer. Here is the question almost nobody can answer: \textbf{did it
actually \emph{need} the corpus?} Maybe it read your rule and used it. Or maybe it already knew
the answer from its own training and your corpus added nothing---in which case the moment your
rules change, or the questions shift, the whole thing quietly breaks and you won't see it
coming. That is the \textbf{necessity} question. There is a twin question: is the corpus
\textbf{sufficient}---does it actually contain what's needed to do the job? Normally you'd find
out by running an expensive study with real people. We wanted a cheap way to answer both
questions \emph{before} involving anyone.

\section*{Idea 1: a test with a right answer (the ``Transfer Instrument'')}

Multiple-choice quizzes are a weak way to measure this, because a fluent-sounding answer can
still be wrong in practice. So instead we drop the learner into a \textbf{realistic task that
has a measurable best move}: steering a ship to avoid a collision. For any situation, a
standard piece of autopilot software computes a strong ``reference'' move. We then score the
learner by \textbf{how far its decision lands from that reference move}. Small distance = good;
large distance = bad. We call that distance \emph{regret}. Unlike a quiz score, regret is
smooth---it can tell ``slightly off'' from ``about to crash.''

We also build the rulebook so that \textbf{each rule controls exactly one part of the score}.
Remove the ``turn to starboard'' rule and \emph{only} the ``did you turn the right way'' part of
the score gets worse; nothing else moves. We check this actually holds---it's not automatic. That
``one rule $\to$ one number'' property is what lets us test rules one at a time.

\section*{Idea 2: run the test backwards to measure NECESSITY}

Here is the core move: \textbf{take a corpus item out and re-run the test.} If the score gets
worse, the learner was genuinely relying on that item---it \emph{needed} it. If the score
\emph{doesn't budge}, the learner already knew the answer and never needed your corpus. That
gap---how much performance falls when we remove an item---\emph{is} the necessity of that item.
Same instrument that grades learning, run in reverse, now audits the corpus.

\section*{The headline: proving necessity is real, by \emph{building} the ground truth}

Here's the honest soft spot, and how we close it---this is the biggest change in the whole
project. When we remove an item and the score doesn't move, we \emph{conclude} the model already
knew the answer. But how do we \emph{know} it already knew it? We could be fooling ourselves. The
only airtight way to prove necessity is real, and not confounded by what the model happened to
already know, is to \textbf{control what the model knows}---and watch necessity move in lockstep.

So we do exactly that: we \textbf{teach a fact directly into a model's weights} and watch its
necessity for that fact fall. Before teaching, the model must read the corpus to answer, so
necessity is high. After teaching, it knows the fact on its own, so necessity should drop to
near zero. If it does, the instrument is measuring what we claim. We ran this two ways, in two
different domains:

\begin{itemize}[leftmargin=1.4em]
  \item \textbf{Ship hazard (a simulator domain).} We invent a hidden charted hazard---a danger
  no model could already know---and teach it into the weights. Necessity collapses from large
  ($\approx$667 in our regret units) to essentially zero ($\approx$0.2). We validated the two
  ends on a \emph{real} model, and we built the ``knowledge gradient'' two independent ways
  (turning a training knob up, and saving the model at successive training checkpoints); the two
  methods \textbf{agree}, which rules out the result being a quirk of how we trained. But it's a
  \emph{step}, not a smooth slope---because this is \emph{one discrete fact}. The model either
  knows it or it doesn't; there's no ``halfway.''
  \item \textbf{Fictional-facts Q\&A (NO simulator---just an answer checker).} Here there is no
  ship, no autopilot, no hand-built barrier: just made-up facts about fictional places that no
  model could possibly know, and a plain checker that asks ``is the answer right?'' We teach those
  facts into a real 7B model, a bit at a time, and necessity falls \textbf{smoothly and steadily:
  1.00 $\to$ 0.93 $\to$ 0.29 $\to$ 0.03} as the model soaks the facts up. It's smooth \emph{because
  there are many independent facts}---you can learn a quarter of them, then half, then almost all.
  This is the clean ``dose-response'' curve: more knowledge in the weights, less need for the
  corpus, every step of the way.
\end{itemize}

The two-domains point matters, and one of them uses \textbf{no simulator at all}. That's what
shows the method isn't a lucky trick of one hand-built ship simulator---it works on any task with
a checkable right answer, and it's calibrated against knowledge we \emph{put in ourselves}, so
``necessity'' provably means ``needed the corpus,'' not ``happened to correlate with something.''

\section*{What this is \emph{for}: two products}

\textbf{Product 1 --- Corpus diagnosis.} Run the test backwards over every rule or fact and
\textbf{rank them by how much the learner relies on each one}. That already tells you which items
are carrying weight and which are dead weight. The sharp part is splitting ``the corpus didn't
help'' into two \emph{opposite} cases that look identical from the outside:
\begin{itemize}[leftmargin=1.4em]
  \item \textbf{Redundant} --- the model already knew it, so removing it changes nothing.
  \emph{Verdict: delete it, it's clutter.}
  \item \textbf{Unusable} --- the model fails the task \emph{even with the item right in front of
  it}. The item has value; this model just can't act on it. \emph{Verdict: the problem is the
  model, not the corpus.}
\end{itemize}
Removing-and-watching (necessity alone) can't tell these apart---both show ``no change.'' What
separates them is the score \emph{with} the item present: near-perfect $\Rightarrow$ redundant;
still failing $\Rightarrow$ unusable. We saw both live on the ship hazard: Claude Haiku and Sonnet
read it \emph{corpus-bound} (they use it), while Llama-70B grounds anyway and \emph{ignores} it
(regret stays $\approx$1207 with the rule present)---the textbook ``unusable'' case.

\textbf{Product 2 --- Corpus sufficiency, and the FALSE SUFFICIENCY alarm.} This is the
diagnostic ordinary RAG scoring simply cannot produce. Imagine a corpus that looks totally
sufficient: the model answers everything correctly, every check is green. Our instrument can tell
you the accuracy is a mirage---the model is \textbf{coasting on what it already knew}, and the
corpus is contributing nothing (necessity $\approx$0 everywhere). We call that \textbf{false
sufficiency}, and it's a warning: this system will fail the instant the questions drift outside
what the model happens to have memorized. \emph{Honest caveat:} sufficiency is always ``sufficient
for \emph{these} questions.'' We can only test the questions we ask, so we always report how many
we asked---and the false-sufficiency alarm is precisely the measurable warning that a
looks-fine-on-this-question-set corpus is fragile.

\section*{How this differs from existing tools (RAGAS, ContextCite)}

Tools like RAGAS and ContextCite are good, and close by, so it's worth being precise. They ask
whether an answer is \textbf{supported by} or \textbf{influenced by} the provided context---did
these words come from that passage? We ask a different question: did the model \textbf{need} the
context at all? And crucially, we \textbf{calibrate} that against a known baseline by teaching
facts in ourselves, so a fact the model already knew is never miscounted as one it relied on.
That calibration---controlling what's in the weights and checking necessity tracks it---is the
thing those tools structurally lack. It's the difference between ``the answer matches the
passage'' and ``the answer would have been wrong without the passage.''

\section*{Additional breadth (clearly separate from the central story)}

Everything above is the \textbf{central story}: measure necessity, prove it's real by construction
in two domains, and turn it into diagnosis and sufficiency products. The following are supporting
breadth---they widen the evidence, but the paper does not lean on them.

\begin{itemize}[leftmargin=1.4em]
  \item \textbf{Cross-model discrimination.} We run the instrument across many different LLMs
  (Claude Haiku/Sonnet, Llama, Nova, and others) and it tells them apart: every model reads the
  \emph{textbook} ship rules as redundant (they all already know the COLREGs), while the invented
  hidden hazard splits them---corpus-bound vs.\ unusable vs.\ inconclusive. Different models,
  genuinely different verdicts.
  \item \textbf{A second simulator domain.} Besides the continuous steering task, we built a
  \emph{discrete, step-by-step} task---a roadside tire change---and showed the same machinery
  works there too (each skill controls exactly one part of the score, and competence and
  performance move together). So the ``one rule, one number'' property isn't special to ships.
  \item \textbf{The unlearning ``says $\ne$ does'' sidelight.} An earlier version of this project
  made machine-\emph{unlearning}---surgically erasing a rule from the model's brain---the
  load-bearing proof. It isn't anymore; the teaching-in dose-response above is the primary
  weight-level validation. But the unlearning result is a lovely illustration in its own right:
  when we \emph{gently} erase ``turn to starboard,'' the model stops \emph{saying} the rule (it
  quits citing ``Rule 14,'' and every text-based ``did it forget?'' score reports success---probe
  dropped from 0.43 to 0.27), yet on the actual steering test the ship \textbf{still turns to
  starboard, unchanged} ($\Delta = 0.000$). The model changed what it \emph{says}, not what it
  \emph{does}---and our instrument measures the \emph{decision}, so it wasn't fooled. A supporting
  sidelight, not the keystone.
  \item \textbf{Honest limitations.} (1)~Sufficiency is only ever ``sufficient for the questions we
  asked''---a bounded question set (see false sufficiency above). (2)~Everything here is
  \textbf{simulation and mechanism evidence, not a human study}. A real trial with real learners is
  written down as the next step, not skipped. (3)~We designed our rulebooks so each rule controls
  one thing; real corpora are messier, and when rules overlap the diagnosis gets coarser (it points
  at a \emph{group} of items, not one).
\end{itemize}

\section*{What we are \emph{not} claiming (the honest part)}

\begin{itemize}[leftmargin=1.4em]
  \item \textbf{This is a test bench, not proof that students learn more.} Everything is
  simulation. The real human study is the next step, not a result in hand.
  \item \textbf{We are not claiming a new AI tutor, or a new collision-avoidance method.} We
  borrow mature ship-autopilot software as a \emph{measuring stick}. The measurement method is the
  contribution.
  \item \textbf{We are not claiming the corpus is ``sufficient'' in general}---only for the
  specific questions we probed. The false-sufficiency alarm is our honest handle on that limit.
\end{itemize}

\section*{Does the logic survive a plain retelling? (the checks)}

\begin{enumerate}[leftmargin=1.4em]
  \item The score is a \emph{distance from an expert's move}, so ``better'' is a measurable
  quantity, not a vibe. \hfill[measurement]
  \item Removing a corpus item and watching the score fall measures \emph{necessity}---whether the
  learner truly needed it. \hfill[the core method]
  \item We prove necessity is real by \emph{teaching a fact into the weights} and watching necessity
  fall---in two domains, one with no simulator, giving a smooth 1.00 $\to$ 0.93 $\to$ 0.29 $\to$
  0.03 curve. \hfill[known-groups validation, the headline]
  \item Read backwards, the same test ranks corpus items and splits ``didn't help'' into
  \emph{redundant} (delete it) vs \emph{unusable} (fix the model). \hfill[corpus diagnosis]
  \item A corpus can look perfectly sufficient while the model coasts on priors---necessity
  $\approx$0 everywhere is the measurable \emph{false-sufficiency} alarm ordinary RAG scoring can't
  raise. \hfill[sufficiency]
  \item We measure whether the model \emph{needed} the context (calibrated against a known
  baseline), where RAGAS/ContextCite measure whether it was \emph{supported/influenced} by it.
  \hfill[the delta vs prior tools]
\end{enumerate}

If any of these feels like a leap, that's exactly where the technical paper needs to work
harder---and where the remaining work (more rules, more models, a real charted-hazard leg, and
eventually the human trial) would turn ``plausible'' into ``demonstrated.''

\end{document}
